\section{The prescribed cyclotomic prism and the full Wach-existence return}
\label{sec:wach-prism-v056}

\status{This chapter constructs the actual prism required by
\texttt{exists\_\allowbreak wach\_\allowbreak prism\_\allowbreak over\_\allowbreak integers}, including its original
scalar map and both coefficient-coordinate assertions. It uses the
already reconstructed integral period rings and invariant matrix
descent. The final return proves the full original Wach-existence
implication under its original tied-period and crystalline-realization
inputs. Later faithfulness, cohomology and quantitative-comparison
producers remain separate.}

The construction has a particularly concrete underlying ring. Let
$p$ be the original prime, $K\subset\overline{\mathbb Q}_p$ the
original finite extension, $K_0$ its maximal unramified subfield,
and put
\[
 W_0=\mathcal O_{K_0},\qquad O=\mathcal O_{LL},\qquad
 \mathfrak S=W_0[[T]],\qquad A_{\mathrm{inf}}=W(\mathcal O_C^\flat).
 \tag{WPR.1}
\]
For the prism endpoint keep the supplied
$P,AI,\mathrm{Good}(K,s,LL),\Gamma,R,I$ and
$\mathrm{IoW}(I)$. In particular $s\geq1$ and $K\subset K_s$,
where $K_s=K_0(\mu_{p^s})$. No new field or cyclotomic parameter
is substituted for these data. The intermediate statements keep
their own printed scopes: in particular WPR11--WPR15 require
no $\mathrm{Good}$ hypothesis, and WPR15 holds for every $s\geq0$. Write $n=[K_0:\mathbb Q_p]$ and
$e_i:K_0\hookrightarrow LL$, $i\in\mathbb Z/n\mathbb Z$, for
the supplied embeddings. The actual coefficient map is
\[
 a:\mathfrak S\longrightarrow\prod_i O[[T]],\qquad
 a(f)_i=e_i(f),
 \tag{WPR.2}
\]
with each $e_i$ applied coefficientwise. The target is the original
$\mathrm{LetA}$, not an independently chosen product.

WIS1--WIS5 supply the complete integral coefficient rings, actual
Witt coordinates and inclusions. IFD9 supplies the unramified
valuation facts and finite cyclotomic integral calculation. LS4
supplies the actual Fontaine map $\theta:A_{\mathrm{inf}}\to\mathcal O_C$
and its principal kernel. These inputs retain their original maps.
We spell out the additional algebra and every prism field below.

\subsection{The unramified constants and the original power-series action}

% BEGIN WPR1
\paragraph{WPR1: unramified integral coordinates and the coefficient Frobenius.}
\label{claim:WPR1}
The ring $W_0$ is a complete DVR with maximal ideal $pW_0$.
It has a finite $\mathbb Z_p$-basis consisting of roots of unity
of order prime to $p$. Consequently every $w\in W_0$ is a
finite $\mathbb Z_p$-linear combination of such roots. The
original coefficient endomorphism $\sigma:W_0\to W_0$ satisfies
\[
                  \sigma(w)=w^p+p b_w\qquad(b_w\in W_0).
 \tag{WPR.3}
\]

\begin{proof}
The complete-DVR and uniformizer assertions are WIS1 and IFD9,
on this same finite unramified subfield. Let its residue field be
$k_0=\mathbb F_{p^f}$. Choose an $\mathbb F_p$-basis
$\bar\omega_1,\ldots,\bar\omega_f$ of $k_0$; each basis element
is nonzero. IFD9 constructs its unique lift $\omega_i$ satisfying
$\omega_i^{p^f-1}=1$. In particular its order is prime to $p$.

For any $w=w_0\in W_0$, choose digits $d_{i,r}\in\{0,\ldots,p-1\}$
successively so that
\[
 w_r=\sum_i d_{i,r}\omega_i+p w_{r+1}.
\]
The residue-basis property permits this step, and divisibility
by $p$ follows from the actual maximal ideal. Thus
\[
 w=\sum_i\omega_i\sum_{r<N}p^r d_{i,r}+p^Nw_N.
\]
All $w_N$ are integral, so the remainder tends to zero. Completeness
gives $w=\sum_i c_i\omega_i$ with
$c_i=\sum_{r\geq0}p^r d_{i,r}\in\mathbb Z_p$. If such a sum is
zero, reducing modulo $p$ makes every $c_i$ divisible by $p$.
Cancel $p$ and repeat; $p$-adic separation of $\mathbb Z_p$ gives
$c_i=0$. This proves the claimed basis and the original finite
root-of-unity spanning statement.

The map $\sigma$ is the restriction of the supplied
$P.\mathrm{PD}(K).\sigma$, transported to the integral ring by
the original $\mathrm{embW}$. It preserves integral elements.
The defining residue-Frobenius property of that datum says
$|\sigma(w)-w^p|<1$ for integral $w$. Since the value group of
$K_0$ is $p^{\mathbb Z}$, every such difference is in $pW_0$.
This gives (WPR.3), including the actual coefficient $b_w$.
\end{proof}
% END WPR1

% BEGIN WPR2
\paragraph{WPR2: the integral cyclotomic character at its full original scope.}
\label{claim:WPR2}
There is a continuous character
$\chi:G_{\mathbb Q_p}\to\mathbb Z_p^\times$ whose image in
$\mathbb Q_p^\times$ is the supplied $P.\chi_Q$. For every finite
$K$, every $r\geq0$ with $K\subset K_r$, and every
$u\in\mathbb Z_p^\times$ satisfying $p^r\mid u-1$, there is
$g\in G_K$ with $\chi(g)=u$. This is the full original
\texttt{exists\_integral\_cyclotomicChar} conclusion.

\begin{proof}
The supplied cyclotomic predicate has the following concrete
content: if $\zeta^{p^m}=1$, there is a natural $b$ such that
$g\zeta=\zeta^b$ and $|P.\chi_Q(g)-b|\leq p^{-m}$.
For $m=0$ this gives $|P.\chi_Q(g)|\leq1$. Applying it to
$g^{-1}$ and using the character identity gives equality.
Thus $P.\chi_Q(g)$ and its inverse are integral, defining the
unique lift $\chi(g)\in\mathbb Z_p^\times$. Its multiplication
law follows by injectivity into $\mathbb Q_p$.

For a primitive $p^m$-th root, the same predicate says that
$\chi(g)\bmod p^m$ is exactly its exponent under $g$. The
stabilizer of that finite extension is open in the original
Galois topology. Therefore each finite congruence of $\chi$ is
locally constant, proving continuity into the inverse limit
$\mathbb Z_p^\times$. This also proves
\[
 g\text{ fixes }\mu_{p^m}
 \quad\Longleftrightarrow\quad \chi(g)\equiv1\pmod{p^m}.
 \tag{WPR.4}
\]

For surjectivity with the stated restriction, the degree and
noncancellation calculation in IFD9 gives
$[K_0(\mu_{p^m}):K_0]=p^{m-1}(p-1)$ for $m\geq1$.
Hence every unit exponent modulo $p^m$ defines an automorphism
of this field over $K_0$. The exponents supplied by $u$ are
compatible as $m$ varies, so they define an automorphism of
$K_\infty=\bigcup_mK_0(\mu_{p^m})$. Extend it to the algebraic
closure: the usual extension of an embedding along algebraic
elements, or its maximal-extension proof, produces an automorphism
of the algebraic closure because the target is algebraically closed
and the extension is normal. It fixes $K_0$, and (WPR.4) shows
that it fixes $K_r$ when $u\equiv1\pmod{p^r}$. It therefore fixes
the supplied $K\subset K_r$. Its character agrees with $u$
modulo every $p^m$, and separatedness gives $\chi(g)=u$.
For $r=0$ the assertion just uses $K\subset K_0$ and the same
construction. No odd-prime restriction was used.
\end{proof}
% END WPR2

% BEGIN WPR3
\paragraph{WPR3: the full cyclotomic power-series action.}
\label{claim:WPR3}
For every commutative $\mathbb Z_p$-algebra $D$ and $c\in\mathbb Z_p$,
put
\[
 B_c(T)=\sum_{j\geq0}\binom cjT^j,\qquad
 \beta_c(T)=B_c(T)-1,\qquad
 \gamma_c(f)=f(\beta_c(T)),\qquad \phi_0=\gamma_p.
 \tag{WPR.5}
\]
Coefficients are first formed in $\mathbb Z_p$ and then mapped
to $D$. These are the source's actual $\mathrm{cyc}$ and
$\mathrm{act}$. They satisfy all its action laws:
\[
 \begin{gathered}
 \gamma_c(T)=\beta_c(T),\quad
 \gamma_c\gamma_d=\gamma_{cd},\quad \gamma_1=1,\quad
 \beta_k(T)=(1+T)^k-1\quad(k\in\mathbb N),\\
 [T]\beta_c=c,\qquad
 \gamma_c\big((1+T)^k-1\big)
       =\beta_c\big((1+T)^k-1\big).
 \end{gathered}                                                     \tag{WPR.6}
\]
The constant coefficient is preserved. Coefficientwise base change
along every $\mathbb Z_p$-algebra map commutes with the action.
If $c$ is a unit, $\gamma_c$ is an automorphism with inverse
$\gamma_{c^{-1}}$.

\begin{proof}
For fixed $j$, the polynomial $c(c-1)\cdots(c-j+1)/j!$ is
continuous as a $\mathbb Q_p$-valued function. On natural integers
its values are integers. Naturals are dense in $\mathbb Z_p$ and
$\mathbb Z_p$ is closed, so these values lie in $\mathbb Z_p$
for every $c\in\mathbb Z_p$. This justifies (WPR.5).
The series $\beta_c$ has zero constant term, so every coefficient
of substitution is a finite sum. Substitution is therefore an
honest algebra endomorphism over an arbitrary $D$.

In $\mathbb Z_p[[T]]$, each coefficient of
$B_{c+d}=B_cB_d$ and $\beta_c(\beta_d(T))=\beta_{cd}(T)$ is a
continuous function of $c,d$. Both identities hold for natural
integers by the ordinary binomial theorem, hence for all $c,d$
by density. Transport them coefficientwise to $D$. They give
composition, the identity and all the displayed formulas. The
last formula uses the commutation of $\gamma_c$ with $\gamma_k$.
Constant terms are preserved because substitution has zero
constant term. Every coefficient expression is a finite algebraic
expression in the original coefficients and the images of
$\binom cj$, so it commutes with the indicated base change.
Composition with $c^{-1}$ proves the unit assertion.
\end{proof}
% END WPR3

% BEGIN WPR4
\paragraph{WPR4: formal-series homomorphisms and their differences.}
\label{claim:WPR4}
Let $f,g:A[[T]]\to D[[T]]$ be ring maps agreeing on constants
and on $T$, with $f(T)\in TD[[T]]$. Then $f=g$.
In particular an $A$-algebra endomorphism $f$ with $f(T)\in TA[[T]]$
preserves constant coefficients and satisfies
\[
                         f(T)-T\mid f(h)-h
                  \qquad(h\in A[[T]]).
 \tag{WPR.7}
\]
No continuity hypothesis is added to these original algebraic statements.

\begin{proof}
Write $h=h_{<N}+T^Nh_N$, with $h_{<N}$ the degree-$<N$
polynomial truncation. The two maps agree on that polynomial,
so $f(h)-g(h)$ is divisible by $f(T)^N$, hence by $T^N$.
This holds for every $N$, and coefficientwise separation gives
equality. For an $A$-algebra endomorphism, comparison with
substitution at $y=f(T)$ proves $f(h)=h(y)$. The constant term
statement follows at once. Explicitly, a quotient for (WPR.7) is
\[
 \sum_{n\geq1}h_n\sum_{j=0}^{n-1}y^jT^{n-1-j},
              \qquad h=\sum_{n\geq0}h_nT^n.
\]
Its degree-$m$ coefficient has contributions only from $n\leq m+1$,
because $y\in(T)$. It is a well-defined formal series, and the
finite geometric-sum identity proves the claimed product equality
coefficient by coefficient.
\end{proof}
% END WPR4

% BEGIN WPR5
\paragraph{WPR5: uniqueness of the commuting cyclotomic action.}
\label{claim:WPR5}
Let $D$ be a domain with $p\neq0$ and $p$ not a unit.
Suppose $f,g$ are $D$-algebra endomorphisms of $D[[T]]$,
their images of $T$ have zero constant term and the same linear
coefficient, and both commute with the endomorphism
$\phi_0(T)=(1+T)^p-1$. Then $f=g$.
Consequently the full original \texttt{act\_unique} holds:
on its complete DVR, the prescribed Frobenius is $\phi_0$,
and an endomorphism commuting with it, of zero constant term and
linear coefficient $c$, is $\gamma_c$.

\begin{proof}
Put $F=f(T)$ and $G=g(T)$. If $F-G\neq0$, let $k\geq2$ be
the first degree where its coefficient $b$ is nonzero.
Commutation gives
\[
 (F-G)((1+T)^p-1)=(1+F)^p-(1+G)^p.
\]
The degree-$k$ coefficient on the left is $p^k b$; on the right
it is $p b$. Every higher power of $F$ or $G$ contributes to
their difference only in degrees at least $k+1$. Thus
$p(p^{k-1}-1)b=0$. None of these factors is zero: if
$p^{k-1}=1$, then $p$ would be a unit. This contradicts the
domain hypothesis. Therefore $F=G$, and WPR4 gives $f=g$.

In the original uniqueness statement, WPR4 first identifies the
given $\phi$ from its value on $T$. WPR3 shows that $\gamma_c$
has the required linear coefficient and commutes with that $\phi$.
Apply the preceding argument. This proves the full original
complete-DVR specialization and, in fact, the broader uniqueness
statement just printed.
\end{proof}
% END WPR5

% BEGIN WPR6
\paragraph{WPR6: identification of the supplied series operators.}
\label{claim:WPR6}
For the original $R:\mathrm{WSer}(\Gamma)$ and the character
of WPR2,
\[
 R.\varphi_0=\phi_0,
 \qquad R.\gamma_0(\mathrm{res}_\Gamma g)=\gamma_{\chi(g)}
                         \quad(g\in G_K).
 \tag{WPR.8}
\]
These are equalities on the actual coefficient ring $O[[T]]$.

\begin{proof}
The source fields $R.\varphi_0(T)=Tq_0$ and the defining
formula for $q_0$ give $R.\varphi_0(T)=(1+T)^p-1$;
the prime of $P.\mathrm{PD}(K)$ is the original $p$.
The supplied $R.\gamma_0(\mathrm{res}_\Gamma g)$ commutes with
this Frobenius, sends $T$ to a series with zero constant term,
and has linear coefficient $\chi_\Gamma(\mathrm{res}_\Gamma g)$.
The original character compatibility identifies this coefficient
with the image of $P.\chi_Q(g)$, hence of $\chi(g)$ in $O$.
WIS1 gives $p\neq0$ and $p$ nonunit in this complete DVR.
WPR5 proves both equalities. This argument quantifies over every
original $g\in G_K$ and does not replace the supplied $R$.
\end{proof}
% END WPR6

% BEGIN WPR7
\paragraph{WPR7: all cyclotomic parameter and ideal-generator formulas.}
\label{claim:WPR7}
For every commutative $\mathbb Z_p$-algebra $D$ and $s\geq1$, put
\[
 U_s=(1+T)^{p^{s-1}}-1,\quad
 Q_s=\sum_{j=0}^{p-1}(1+U_s)^j,\quad d_s=p^{s-1}(p-1).
 \tag{WPR.9}
\]
Then the full original \texttt{XX\_QQ\_properties} holds:
\[
 \begin{gathered}
 \phi_0(U_s)=U_sQ_s=(1+T)^{p^s}-1,
 \qquad \phi_0^{s-1}(T)=U_s,\quad Q_s(0)=p,\\
 Q_s=T^{d_s}+p b,\qquad Q_s=p+U_s b'
                    \quad\text{for some }b,b'\in D[[T]],\\
 \gamma_c(U_s)=U_s(c+U_s b_c),\quad
 \gamma_c(Q_s)=Q_s(1+U_s b'_c)\quad(c\in\mathbb Z_p^\times),\\
 u=1+p^s w\quad\Longrightarrow\quad
         \gamma_u(f)-f\in U_sQ_sD[[T]]\quad(f\in D[[T]]).
 \end{gathered}                                                     \tag{WPR.10}
\]

\begin{proof}
The geometric-sum identity gives the first product, and iteration
of $\phi_0(T)=(1+T)^p-1$ gives the second. At $T=0$ each of the
$p$ summands is one. The polynomial
$\sum_{j<p}(1+Z)^j$ is congruent to $Z^{p-1}$ modulo $p$:
multiply by $Z$ and use $(1+Z)^p-1\equiv Z^p$ in
$\mathbb F_p[Z]$, then cancel $Z$ in that polynomial ring.
Lifting the integral coefficient identity to any $D$, and using
$U_s\equiv T^{p^{s-1}}\pmod p$, gives
$Q_s=T^{d_s}+pb$. Subtracting the value at $U_s=0$ gives
$Q_s=p+U_sb'$ directly.

The zero and linear coefficients in WPR3 give
$\beta_c(T)=T(c+Tb_c)$. Commuting with $\gamma_{p^{s-1}}$
gives the formula for $U_s$. For $Q_1$ write $b=b_c$ and
$D_c=c+Tb$. It is a unit since $c$ is. Commutation applied to
$\phi_0(T)=TQ_1$ gives, after cancellation of the regular
formal variable $T$,
\[
 D_c\gamma_c(Q_1)=Q_1\{c+TQ_1\phi_0(b)\}.
\]
Thus
$\gamma_c(Q_1)=Q_1[1+TD_c^{-1}(Q_1\phi_0(b)-b)]$.
Applying $\gamma_{p^{s-1}}$ proves the asserted formula for
$Q_s$. Cancellation here is by $T$ in a power-series ring,
which is valid over every coefficient ring $D$.

Finally set $V=U_sQ_s=(1+T)^{p^s}-1$. The additive and
multiplicative binomial identities in WPR3 show
\[
 \gamma_{1+p^s w}(T)-T=(1+T)\beta_w(V).
\]
The right side is divisible by $V$, because $\beta_w$ has
zero constant term. WPR4 then proves the last assertion for
every formal series $f$. All arguments commute with the given
coefficient map, including rings with $p$-torsion.
\end{proof}
% END WPR7

\subsection{Division, completeness and cyclotomic specialization}

% BEGIN WPR8
\paragraph{WPR8: full distinguished division and the original kernel criterion.}
\label{claim:WPR8}
Let $D$ be a complete separated local ring for its maximal ideal,
let $p\in\mathfrak m_D$, and let
$Q=T^d+pb\in D[[T]]$. Every $h\in D[[T]]$ has a unique
expression
\[
                     h=Qq+r,\qquad r\in D[T],\quad\deg r<d.
 \tag{WPR.11}
\]
For $d=0$ this means $r=0$. If a ring map $f:D[[T]]\to C$
satisfies $f(Q)=0$ and is injective on polynomials of degree
less than $d$, then
\[
                              \ker f=(Q).
 \tag{WPR.12}
\]
This proves the full original \texttt{ker\_eq\_span}, with no
Noetherian or continuity assumption added.

\begin{proof}
For $d=0$, $Q=1+pb$ has unit constant coefficient, so it is a
unit in $D[[T]]$ and the statement follows. Suppose $d>0$.
Write $H_d(h)$ for the tail after removing its first $d$
coefficients, so $h=r_0+T^dH_d(h)$ with $\deg r_0<d$.
The required quotient must solve
\[
       q+pL(q)=H_d(h),\qquad L(q)=H_d(bq).
\]
It is given coefficientwise by
\[
                       q=\sum_{j\geq0}(-p)^jL^j(H_d(h)).
\]
For each coefficient, the $j$th term lies in
$\mathfrak m_D^j$, so completeness defines the sum.
Each coefficient of $L(q)$ uses only finitely many coefficients
of $q$ and $b$; hence $L$ commutes with these limits.
Telescoping gives $q+pL(q)=H_d(h)$. Therefore the tail of
$h-Qq$ vanishes, yielding the required polynomial $r$.

If $q+pL(q)=0$, iteration gives
$q=(-p)^NL^N(q)$ for every $N$. Each coefficient belongs to
every $\mathfrak m_D^N$, so separation makes $q=0$.
The quotient and then the remainder are unique. This proof uses
the given maximal-ideal completeness, not a tacit assumption that
every principal ideal is closed in that topology.

Apply $f$ to (WPR.11). If $f(h)=0$, then $f(r)=0$ because
$f(Q)=0$. The stated degree-$<d$ injectivity makes $r=0$,
so $h\in(Q)$. The reverse inclusion is immediate. This also
handles the possible zero polynomial without a natural-degree
convention causing an extra case.
\end{proof}
% END WPR8

% BEGIN WPR9
\paragraph{WPR9: completeness for the actual prism ideal.}
\label{claim:WPR9}
Let $D$ be complete and separated for $(p)$, and suppose
$Q=T^d+pb$ with $d>0$. Then $D[[T]]$ is complete and
separated for $(p,Q)$. This is the original
\texttt{isAdicComplete\_span} at its full coefficient-ring scope.

\begin{proof}
First put $J=(p,T)$. A series belongs to $J^N$ exactly when
its coefficient of $T^i$, for $i<N$, belongs to $p^{N-i}D$.
One implication follows from the generators $p^{N-i}T^i$;
for the other, choose the indicated finitely many quotients and
leave the degree-$\geq N$ tail in $T^ND[[T]]$.

A compatible system modulo $J^N$ therefore determines, at every
fixed degree, a compatible system of coefficients modulo all
powers of $p$. Completeness of $D$ supplies a unique coefficient.
The resulting formal series has every prescribed class by the
same membership criterion. Separation follows coefficientwise.
Thus $D[[T]]$ is $J$-adically complete and separated.

Set $I=(p,Q)$. From $T^d=Q-pb$ and $d>0$ we obtain
$J^d\subset I\subset J$. Consequently
$J^{dN}\subset I^N\subset J^N$ for all $N$. These explicit
containments identify the Cauchy systems and separated limits
for the two topologies. They transfer the completeness just
proved to $I$, without finite-generation assumptions on $D$.
\end{proof}
% END WPR9

% BEGIN WPR10
\paragraph{WPR10: the full integral cyclotomic polynomial statement.}
\label{claim:WPR10}
Let $s\geq1$, let $\zeta$ be any primitive $p^s$-th root in
the original algebraic closure, and let
$t\in\mathcal O_C$ be its specified image $\zeta-1$.
Then
\[
 \deg r<d_s,\quad r\in W_0[T],\quad r(t)=0
                       \Longrightarrow r=0.                  \tag{WPR.13}
\]
If $K\subset K_s$, every $x\in\mathcal O_K$ has
$x=r(t)$ in $\mathcal O_C$ for some $r\in W_0[T]$.
All embeddings are the original $\mathrm{intToOC}$ maps.

\begin{proof}
The root-norm calculation in LS2, used in IFD9, applies to any
primitive root and gives $|\zeta-1|=p^{-1/d_s}$. For
$a_i\in K_0$, the nonzero terms $a_i(\zeta-1)^i$,
$0\leq i<d_s$, have distinct valuations modulo $\mathbb Z$,
because the valuation of each $a_i$ is an integer. Their sum
cannot vanish. This proves (WPR.13).

These $d_s$ powers are a basis of $K_s$ over $K_0$: the preceding
independence gives the lower degree bound, and the cyclotomic
polynomial gives the upper bound. Write an integral $x\in K_s$
in this basis. The norm of the sum is its largest term norm.
Since the exponents $i/d_s$ lie in $[0,1)$ and the coefficient
valuations are integral, every coefficient is integral.
Thus $\mathcal O_{K_s}=W_0[\zeta-1]$. Apply this to the given
$x\in\mathcal O_K\subset\mathcal O_{K_s}$. Polynomial evaluation
commutes with the specified inclusions into the algebraic closure
and then its completion; hence the resulting equality is exactly
the original one in $\mathcal O_C$. The case $p=2,s=1$, where
$d_s=1$, is included.
\end{proof}
% END WPR10

\subsection{Descending the prescribed tensor-valued inclusion}

% BEGIN WPR11
\paragraph{WPR11: a normalized linear retraction.}
\label{claim:WPR11}
Let $Z$ be a local ring with maximal ideal $(p)$, let $O'$ be
a free $Z$-algebra, and suppose the image of $p$ is not a unit
in $O'$. There is a $Z$-linear map $r:O'\to Z$ with $r(1)=1$.
For every $Z$-algebra $A$, the map
\[
 j:A\longrightarrow O'\otimes_Z A,\qquad x\longmapsto1\otimes x
 \tag{WPR.14}
\]
is injective. The induced linear map
$\rho=r\otimes1:O'\otimes_Z A\to A$ has
$\rho j=1$ and $\rho(j(a)y)=a\rho(y)$.

\begin{proof}
Choose a free basis of $O'$ over $Z$. At least one coordinate
$c$ of $1$ is a unit. Otherwise all its finitely many nonzero
coordinates belong to $(p)$, so $1=pb$ in $O'$ for some $b$,
contradicting the nonunit hypothesis. Multiply the corresponding
coordinate functional by $c^{-1}$ to obtain $r$.
The balanced tensor map $\rho$ satisfies
$\rho(o\otimes a)=r(o)a$, giving $\rho(1\otimes a)=a$.
This proves injectivity of $j$. On a pure tensor,
$\rho(j(a)(o\otimes b))=r(o)ab=a\rho(o\otimes b)$;
additivity gives the formula for every tensor. No finiteness
of the free basis was needed in the retraction argument.
\end{proof}
% END WPR11

% BEGIN WPR12
\paragraph{WPR12: full formal-series factorization through a pure tensor.}
\label{claim:WPR12}
Let $Z,O',A$ be commutative rings with the indicated $Z$-algebra
structures. Suppose a $Z$-linear $r:O'\to Z$ satisfies $r(1)=1$.
Let $D$ be any commutative ring and
$F:D[[T]]\to O'\otimes_Z A$ a ring map such that
\[
 F(T)=j(\pi),\qquad F(D)\subset j(A).
\]
Assume only that $O'\otimes_Z A$ is separated for the ideal
$(j(\pi))$. Then there is a unique ring map
$\iota:D[[T]]\to A$ with $F=j\iota$.
This proves the full original \texttt{exists\_factor} conclusion.

\begin{proof}
Use $\rho$ and the injectivity of $j$ from the tensor calculation
in WPR11; that calculation uses only the supplied $r(1)=1$.
Every polynomial has image in $j(A)$, since its constants and
variable do. For any $h$ and any $N$, write
$h=h_{<N}+T^Nh_N$. Then
\[
 F(h)=j(a_N)+j(\pi)^NF(h_N)
\]
for some $a_N\in A$. The identity
$\rho(j(a)y)=a\rho(y)$ gives
\[
 j\rho(F(h))-F(h)
   =j(\pi)^N\{j\rho(F(h_N))-F(h_N)\}.
\]
The left side belongs to every power of $(j(\pi))$ and is
therefore zero. Thus $F(h)\in j(A)$ for every $h$.
Define $\iota(h)$ to be its unique preimage. Injectivity of $j$
and the ring identities for $F$ prove all ring identities for
$\iota$, and also uniqueness. This argument uses formal
truncations and separatedness; the map $F$ need not be supplied
with any topological structure or continuity property.
\end{proof}
% END WPR12

% BEGIN WPR13
\paragraph{WPR13: separation for every permitted Witt parameter.}
\label{claim:WPR13}
Let $e\in\mathcal O_C^\flat$ with $e-1$ a nonunit and put
$\mu=[e]-1\in A_{\mathrm{inf}}$. Then
\[
 \bigcap_{N\geq0}\mu^NA_{\mathrm{inf}}=0,
 \qquad
 \bigcap_{N\geq0}(1\otimes\mu)^N(O\otimes_{\mathbb Z_p}A_{\mathrm{inf}})=0.
 \tag{WPR.15}
\]
In particular this applies to the original $e=\epsilon_s$.

\begin{proof}
If $e=1$, both intersections are zero immediately. Otherwise
$\tau=e-1$ is a nonzero element of norm less than one in the
tilt valuation ring. Hence $\bigcap_N\tau^N\mathcal O_C^\flat=0$.
The image of $\mu$ modulo $p$ is $\tau$.

Every ideal $(\mu^N)$ is $p$-saturated. Indeed, if
$p x=\mu^N y$, reduction modulo $p$ gives $\tau^Ny_0=0$.
The tilt is a domain, so $y_0=0$ and WIS2 gives $y=pz$.
Cancel $p$ in $A_{\mathrm{inf}}$ to obtain $x=\mu^Nz$.
If $x$ belongs to every $(\mu^N)$, its zeroth Witt coordinate
belongs to every $(\tau^N)$ and is zero. Write $x=py$.
Saturation puts $y$ in the same intersection. Repeating gives
$x\in\bigcap_m p^mA_{\mathrm{inf}}=0$.

Choose a finite $\mathbb Z_p$-basis of $O$, available from WIS1.
It identifies $O\otimes A_{\mathrm{inf}}$ with finitely many copies of
$A_{\mathrm{inf}}$ as an $A_{\mathrm{inf}}$-module. Multiplication by $1\otimes\mu$
is coordinatewise multiplication by $\mu$, proving the second
intersection assertion. Finally WIS3 proves that
$\epsilon_s-1$ is a nonunit for every original $s$; alternatively,
its coordinates are nilpotent modulo $p$ by the cyclotomic
power identities. Thus the original parameter is included.
\end{proof}
% END WPR13

% BEGIN WPR14
\paragraph{WPR14: descent of every unramified constant with its theta value.}
\label{claim:WPR14}
For every $w\in W_0$ there is $b_w\in A_{\mathrm{inf}}$ such that
\[
 I.\mathrm{io}(a(Cw))=1\otimes b_w,
             \qquad \theta(b_w)=\mathrm{intToOC}_{K_0}(w).
 \tag{WPR.16}
\]
Here $Cw$ is the constant formal series, and the supplied
$\mathrm{IoW}(I)$ is the only root-compatibility premise used.

\begin{proof}
For $z\in\mathbb Z_p$, every $e_i$ fixes $z$ and
$I.\mathrm{io}$ is an $O$-algebra map. The tensor balancing
identity therefore gives
$I.\mathrm{io}(a(Cz))=1\otimes z$.
The actual theta map agrees on $\mathbb Z_p$ with its usual
embedding into $\mathcal O_C$. This can be checked without a
continuity assumption on a ring map $\psi:\mathbb Z_p\to\mathcal O_C$:
write $z=b_N+p^Ny_N$ with an integer $b_N$ and
$y_N\in\mathbb Z_p$. The images under $\psi$ and the usual map
differ by $p^N$ times an integral element, so their difference
has norm at most $p^{-N}$ for every $N$. They are equal.

Now let $\omega\in W_0$ be a root of unity of order $m$
prime to $p$. Choose the compatible $p$-power roots of $\omega$
inside its finite cyclic group; explicitly their exponents are
$p^{-r}\bmod m$. Their reductions form the original allowed
tilt $\omega^\flat$, with $(\omega^\flat)^\sharp=\omega$.
Apply $\mathrm{IoW}(I)$ to the actual tuple $(e_i\omega)_i$.
It gives
\[
 I.\mathrm{io}(a(C\omega))=1\otimes[\omega^\flat],
       \qquad \theta([\omega^\flat])=\omega.
\]
WPR1 expresses any $w$ as a finite sum
$\sum_i z_i\omega_i$. Both $I.\mathrm{io}\circ a\circ C$ and
$\theta$ are ring maps, and $b\mapsto1\otimes b$ is a ring
map. Taking $b_w=\sum_i z_i[\omega_i^\flat]$ proves both
equalities. The construction uses a finite linear combination;
no additivity of Teichmuller representatives is assumed.
\end{proof}
% END WPR14

% BEGIN WPR15
\paragraph{WPR15: the full original inclusion into one copy of Ainf.}
\label{claim:WPR15}
For every original $P,AI,K,LL$, every $s\geq0$, and every
supplied $\Gamma,R,I,\mathrm{IoW}(I)$, there is a ring map
$\iota:\mathfrak S\to A_{\mathrm{inf}}$ with
\[
 \begin{gathered}
 I.\mathrm{io}(a(f))=1\otimes\iota(f)\quad(f\in\mathfrak S),\\
 \iota(T)=[\epsilon_s]-1,\qquad
 \theta\iota(Cw)=\mathrm{intToOC}_{K_0}(w)\quad(w\in W_0).
 \end{gathered}                                                     \tag{WPR.17}
\]
This supplies the full original
\texttt{exists\_wachInclusion\_powerSeries\_into\_one\_tmul\_ainf}
on every permitted supplied $\Gamma,R,I,\mathrm{IoW}(I)$.

\begin{proof}
Use WPR11 with $Z=\mathbb Z_p$ and $O'=O$; WIS1 gives
freeness and noninvertibility of $p$. Let $F=I.\mathrm{io}\circ a$.
Its value on $T$ is $1\otimes([\epsilon_s]-1)$ by the original
$I.\mathrm{io}_\pi$ field, because $a(T)$ is the tuple of
the original variables. WPR14 gives all its constant values in
the image of $1\otimes A_{\mathrm{inf}}$. WPR13 provides the required
separation. WPR12 now supplies $\iota$ and the first equality.

Injectivity of $b\mapsto1\otimes b$ then gives the asserted
value of $\iota(T)$ and identifies $\iota(Cw)$ with the
particular $b_w$ of WPR14. This proves its theta value as well.
No $\mathrm{Good}$, $s\geq1$ or $K\subset K_s$ assumption
is used in WPR11--WPR15: separation in WPR13 holds at every
original root index, and WPR14 uses only the supplied constant
map and $\mathrm{IoW}$. All original witnesses therefore remain
arbitrary at this broader scope. The proof even applies
to the source's broader \texttt{exists\_io} statement with any
permitted $e$ and coefficient-tuple map satisfying its printed
constant and variable hypotheses.
\end{proof}
% END WPR15

% BEGIN WPR16
\paragraph{WPR16: the actual Frobenius and Galois intertwining maps.}
\label{claim:WPR16}
Define the ring endomorphism
\[
 \varphi_{\mathfrak S}(f)
       =\sigma(f)((1+T)^p-1)
\]
by applying $\sigma$ coefficientwise before substitution. Then
\[
 \begin{gathered}
 a\varphi_{\mathfrak S}=\mathrm{let}\varphi_{R.\varphi_0}\,a,
 \qquad a\gamma_{\chi(g)}
       =\mathrm{let}\gamma_{R.\gamma_0(\mathrm{res}_\Gamma g)}\,a,\\
 \iota\varphi_{\mathfrak S}=\varphi_{A_{\mathrm{inf}}}\iota,
 \qquad \iota\gamma_{\chi(g)}=g\iota
                                \quad(g\in G_K).
 \end{gathered}                                                     \tag{WPR.18}
\]
The coefficient map $a$ is injective.

\begin{proof}
The source embeddings satisfy $e_i\sigma=e_{i-1}$, with the
original cyclic index convention. Applying WPR3's base-change
formula and then WPR6 gives the first identity componentwise:
coefficient Frobenius shifts precisely to component $i-1$,
and the variable substitution is the prescribed $R.\varphi_0$.
For the second identity no shift occurs; each $e_i$ commutes
with $\gamma_{\chi(g)}$, and WPR6 identifies its source operator.
The index set is nonempty since $[K_0:\mathbb Q_p]>0$.
Evaluation at one component and injectivity of its field embedding
show that $a$ is injective coefficient by coefficient.

To prove the last two identities, first apply
$b\mapsto1\otimes b$ and then the actual inclusion
$O\otimes A_{\mathrm{inf}}\hookrightarrow O\otimes W(\operatorname{Frac}
\mathcal O_C^\flat)$. This composite is injective by WPR11
and WIS5. Substitute (WPR.17) and the first two identities in
(WPR.18). The defining $I.\mathrm{ioE}$ Frobenius and Galois
compatibilities, together with its extension of $I.\mathrm{io}$,
make the two resulting expressions equal in that tensor ring.
On a pure tensor $1\otimes b$, the tensor Frobenius and Galois
actions give $1\otimes\varphi(b)$ and $1\otimes g(b)$.
These equalities follow directly from the Witt coordinate maps.

The Galois action specified by $AI$ is the same action used here:
its required compatibility with every inverse-limit coordinate
determines the tilt action uniquely, and Witt functoriality then
determines its action on $A_{\mathrm{inf}}$. This is also the coordinate
equivariance checked in IFD8. Injectivity of the composite map
now proves the last two identities in $A_{\mathrm{inf}}$ itself.
\end{proof}
% END WPR16

% BEGIN WPR17
\paragraph{WPR17: the exact specialization kernel and its base structure.}
\label{claim:WPR17}
Put $f=\theta\iota:\mathfrak S\to\mathcal O_C$.
Then
\[
 \ker f=(Q_s),\qquad
 \mathrm{intToOC}_K(\mathcal O_K)\subset f(\mathfrak S).
 \tag{WPR.19}
\]
There is a ring map
$\mathrm{str}:\mathcal O_K\to\mathfrak S/(Q_s)$ such that,
for every $x\in\mathcal O_K$ and $b\in\mathfrak S$,
\[
 [b]=\mathrm{str}(x)
       \quad\Longleftrightarrow\quad f(b)=\mathrm{intToOC}_K(x).
 \tag{WPR.20}
\]

\begin{proof}
By WPR15 and the exact sharp-root values of the original
$\epsilon_s$, $f(T)=\zeta_{p^s}-1$ and
$f(Cw)=\mathrm{intToOC}_{K_0}(w)$. Thus polynomial evaluation
under $f$ is the prescribed evaluation in WPR10. Moreover,
\[
 f(Q_s)=\sum_{j=0}^{p-1}\zeta_p^j=0.
\]
Indeed multiplying this sum by $\zeta_p-1\neq0$ gives zero,
and $\mathcal O_C$ is a domain. WPR7 gives
$Q_s=T^{d_s}+pb$, while WPR10 gives the required injectivity
on polynomials of degree less than $d_s$. WPR8 therefore gives
the kernel equality. The second conclusion in WPR10, using the
actual $K\subset K_s$, gives the image assertion in (WPR.19).

Here is the general algebra behind (WPR.20). For ring maps
$f:S\to C$, $g:D\to C$, with $\ker f=J$ and
$g(D)\subset f(S)$, the induced map $\bar f:S/J\to C$
is injective and its image contains $g(D)$. Define
$\mathrm{str}(x)$ to be the unique preimage of $g(x)$.
Injectivity proves preservation of addition, multiplication and
one, and gives $[b]=\mathrm{str}(x)$ exactly when $f(b)=g(x)$.
Apply this to the present maps. This is the full original
\texttt{exists\_str} argument, including its exact equivalence.
\end{proof}
% END WPR17

\subsection{All fields of the three actual prisms and their maps}

% BEGIN WPR18
\paragraph{WPR18: delta structures from the actual Frobenius lifts.}
\label{claim:WPR18}
Let $A$ be a $p$-torsion-free commutative ring and let
$\varphi:A\to A$ be a ring endomorphism with
$\varphi(x)\equiv x^p\pmod{pA}$. There is a unique function
$\delta$ with $\varphi(x)=x^p+p\delta(x)$, and it satisfies
\[
 \begin{split}
 \delta(1)&=0,\\
 \delta(xy)&=x^p\delta(y)+y^p\delta(x)+p\delta(x)\delta(y),\\
 \delta(x+y)&=\delta(x)+\delta(y)
       -\sum_{j=1}^{p-1}\frac{\binom pj}{p}x^jy^{p-j}.
 \end{split}                                                       \tag{WPR.21}
\]
A ring map between two such rings that commutes with Frobenius
also commutes with $\delta$ whenever the target is $p$-torsion-free.
The maps $\sigma$ on $W_0$, $\varphi_{\mathfrak S}$ on
$\mathfrak S$ and Witt Frobenius on $A_{\mathrm{inf}}$ have these structures.

\begin{proof}
Divisibility gives existence of $\delta(x)$ and cancellation of
$p$ gives uniqueness. Expand the ring identities for $\varphi$
and cancel $p$ to obtain (WPR.21); the middle binomial
coefficients are divisible by the prime $p$. Applying a ring
map to $\varphi(x)=x^p+p\delta(x)$ and comparing with the
same identity in the target proves functoriality by cancellation.

WPR1 gives the Frobenius congruence for $\sigma$. For any
commutative $\mathbb Z_p$-algebra $D$ with an endomorphism
$\sigma$ lifting coefficient Frobenius, its associated series
map is congruent modulo $p$ to
$f\mapsto\sum a_i^pT^{pi}=f^p$. Substitution is coefficientwise
finite, and $(1+T)^p-1\equiv T^p$. Thus every coefficient of
$\varphi_{D[[T]]}(f)-f^p$ is divisible by $p$, and choosing
their quotients gives a formal series quotient. This proves the
full original \texttt{phiS\_lift}, including coefficients with
$p$-torsion. For the actual $W_0$ and $\mathfrak S$,
$p$-torsion-freeness follows from the domain property.
For $A_{\mathrm{inf}}$, WIS2 proves $p$-torsion-freeness and identifies
reduction modulo $p$ with the tilt; Witt Frobenius becomes its
$p$th-power map. The constructions therefore apply to all three
actual rings.
\end{proof}
% END WPR18

% BEGIN WPR19
\paragraph{WPR19: the actual Fontaine prism over the original integers.}
\label{claim:WPR19}
The original $A_{\mathrm{inf}}$, its Witt-Frobenius delta structure,
ideal $\ker\theta$, and structure map induced by
$\mathcal O_K\hookrightarrow\mathcal O_C$ form the source's
bounded prism $\mathrm{ainfObj}(K)$. Its Galois endomorphism
for $g\in G_K$ is the original action on Witt coordinates.

\begin{proof}
Use the cyclotomic generator $\xi$ of LS4. That proof gives
$\ker\theta=(\xi)$ and shows that $\xi$ is a nonzerodivisor.
Multiplication by $\xi$ identifies its ideal with the free
rank-one module $A_{\mathrm{inf}}$, proving the required invertibility
of the ideal as a module. WPR18 supplies the delta structure.
The same cyclotomic formula gives
\[
 \theta\varphi(\xi)=p,
 \qquad p=\{p-\varphi(\xi)\}+\varphi(\xi)
                  \in(\xi)+\varphi((\xi)).                 \tag{WPR.22}
\]
The quotient is $\mathcal O_C$ by the actual surjective theta
map. It is $p$-torsion-free, so its bounded-torsion exponent may
be taken to be zero.

We verify the mixed completeness field. Write $S=\mathcal O_C^\flat$
and $\alpha=\xi\bmod p=(\epsilon_1-1)^{p-1}$.
WIS22 shows that $S$ is complete and separated for
$(\epsilon_1-1)$ and hence for $(\alpha)$, since a positive
power gives cofinal ideal powers. LS4 proves that every
$(\xi^m)$ in $A_{\mathrm{inf}}$ is $p$-saturated. Iterating saturation
gives
\[
             (\xi^m)\cap p^NA_{\mathrm{inf}}=p^N(\xi^m).             \tag{WPR.23}
\]
We claim $A_{\mathrm{inf}}/p^N$ is $\xi$-adically complete and separated
for every $N$. For $N=1$ this is the assertion about $S$.
Induct using the exact sequence with left term
$p^{N-1}A_{\mathrm{inf}}/p^NA_{\mathrm{inf}}\simeq S$ and right term
$A_{\mathrm{inf}}/p^{N-1}$. Equation (WPR.23) makes its reductions
modulo $\xi^m$ exact, with left term $S/\alpha^mS$.
Those left-term transition maps are surjective. Compatible
middle lifts may therefore be chosen successively: the defect
of a lift at stage $m+1$ is a class in $S/\alpha^mS$, which
lifts from $S/\alpha^{m+1}S$ and can be subtracted. The same
exact sequence and separatedness of its two outer terms show
uniqueness. This proves the induction assertion.

Now a compatible rectangular system modulo $(p^N,\xi^m)$
first gives an element of $A_{\mathrm{inf}}/p^N$ for each $N$, and
then an element of $A_{\mathrm{inf}}$ by its $p$-adic completeness.
All maps are the original quotient maps, so this element has
the specified rectangular classes and is unique. Finally
\[
 (p,\xi)^{2N}\subset(p^N,\xi^N)\subset(p,\xi)^N
\]
shows that the diagonal rectangular system is cofinal with
the required mixed-adic system. This proves the prism's
completeness and separation fields.

The structure map is the original inclusion into $\mathcal O_C$
followed by the inverse of the theta quotient isomorphism.
Galois acts on the tilt coordinates and hence on Witt vectors,
commutes with Frobenius and with theta, and preserves its kernel.
It therefore preserves $\delta$ by WPR18. Since $g\in G_K$
fixes the original $\mathcal O_K$, it also preserves the
structure map. These are exactly the fields of
$\mathrm{ainfHom}(K,g)$ on the actual $\mathrm{ainfObj}(K)$.
\end{proof}
% END WPR19

% BEGIN WPR20
\paragraph{WPR20: construction of the original series prism.}
\label{claim:WPR20}
The following data form the required bounded prism over
$\mathcal O_K$:
\[
 P_{\mathfrak S}.A=\mathfrak S,
 \qquad P_{\mathfrak S}.I=(Q_s),
 \qquad P_{\mathfrak S}.D.\varphi=\varphi_{\mathfrak S},
 \qquad P_{\mathfrak S}.\mathrm{str}=\mathrm{str}
                     \text{ of WPR17}.                       \tag{WPR.24}
\]
Its ring is local and a domain. Its quotient has no nonzero
$p$-power torsion.

\begin{proof}
The least nonzero coefficient of a product proves the domain
property of $W_0[[T]]$. A formal series is a unit exactly when
its constant coefficient is a unit in $W_0$; construct the
inverse recursively from that coefficient. Its unique maximal
ideal is therefore $(p,T)$, proving locality. WPR18 gives the
delta structure.

The constant coefficient of $Q_s$ is $p\neq0$, so $Q_s$ is
nonzero and regular. Its ideal is free of rank one via
multiplication by $Q_s$, and hence is an invertible module.
WPR1 gives $p$-adic completeness of $W_0$; WPR7 and WPR9
give completeness for $(p,Q_s)$ in $\mathfrak S$.

Write $Q_s=p+U_s b$ using WPR7. Since $\sigma$ fixes its
integral binomial coefficients,
$\varphi_{\mathfrak S}(U_s)=U_sQ_s$. Consequently
\[
 \varphi_{\mathfrak S}(Q_s)
     =p+U_sQ_s\varphi_{\mathfrak S}(b),\qquad
 p=\varphi_{\mathfrak S}(Q_s)
          -Q_sU_s\varphi_{\mathfrak S}(b).
 \tag{WPR.25}
\]
This proves $p\in(Q_s)+\varphi_{\mathfrak S}((Q_s))$.
By WPR17 the quotient embeds into $\mathcal O_C$ under $f$,
so it has no nonzero $p$-power torsion; exponent zero suffices
for the boundedness field. The map in (WPR.20) supplies its
base structure. These are all the fields of the original
$\mathrm{prismS}$, on the specified ring and parameter.
\end{proof}
% END WPR20

% BEGIN WPR21
\paragraph{WPR21: the constant prism and the original evaluation morphism.}
\label{claim:WPR21}
There is a bounded prism $P_W$ on $W_0$, with ideal $(p)$ and
Frobenius $\sigma$, and a prism morphism
\[
 \mathrm{ev}:P_{\mathfrak S}\longrightarrow P_W,
 \qquad \mathrm{ev}(h)=h(0).
 \tag{WPR.26}
\]
This ring map is surjective and has kernel $(T)$.

\begin{proof}
WPR1 and WPR18 supply the complete ring and delta structure.
The nonzero element $p$ is regular in $W_0$, so its ideal is
free of rank one. Completeness for $(p)+(p)$ is exactly
$p$-adic completeness. The prism containment follows already
from $p\in(p)$. Every element of $W_0/(p)$ is killed by $p$,
so the boundedness exponent is one. Give this quotient the
structure map
\[
 \mathcal O_K\xrightarrow{\mathrm{str}}
 \mathfrak S/(Q_s)\xrightarrow{\overline{\mathrm{ev}}}W_0/(p).
\]
The second arrow exists because $Q_s(0)=p$. These data are
exactly the original $\mathrm{prismW}$ with its specified
structure map.

Constant coefficients give
$\mathrm{ev}\,\varphi_{\mathfrak S}=\sigma\,\mathrm{ev}$.
Since $W_0$ has no $p$-torsion, WPR18 gives compatibility with
delta. Evaluation sends $(Q_s)$ into $(p)$, and the displayed
definition makes it compatible with the base structure. This
verifies every field of a prism morphism. The constant series
$C(w)$ maps to $w$, proving surjectivity. A series has zero
constant coefficient exactly when it is $T$ times the series
obtained by shifting all coefficients down by one. Thus the
kernel is precisely $(T)$, with no topological assumption on
the ring map.
\end{proof}
% END WPR21

% BEGIN WPR22
\paragraph{WPR22: the actual period-ring morphism of prisms.}
\label{claim:WPR22}
The ring map $\mathrm{io}$ of WPR15 defines a morphism
\[
 \mathrm{ioH}:P_{\mathfrak S}\longrightarrow\mathrm{ainfObj}(K)
 \tag{WPR.27}
\]
to the prescribed Fontaine prism, with the original pure-tensor
identity $I.\mathrm{io}(a h)=1\otimes\mathrm{io}(h)$.

\begin{proof}
The ring map is the unique factor of the original
$I.\mathrm{io}\circ a$ constructed in WPR15. WPR16 proves
that it commutes with the original Frobenius maps, and WPR18
therefore proves compatibility with delta in the $p$-torsion-free
target. WPR17 says that $f=\theta\,\mathrm{io}$ has kernel
$(Q_s)$, so $\mathrm{io}$ sends the source ideal into
$\ker\theta$, the target ideal.

For $x\in\mathcal O_K$ and $h\in\mathfrak S$, WPR17 gives
\[
 [h]=\mathrm{str}(x)
 \quad\Longleftrightarrow\quad
 \theta\,\mathrm{io}(h)=\mathrm{intToOC}(x).
\]
The same theta criterion defines the structure map of
$\mathrm{ainfObj}(K)$ in WPR19. Hence $[\mathrm{io}(h)]$
is its structure class of $x$ whenever $[h]$ is the source
structure class. This is exactly the remaining base-compatibility
field of a prism morphism. The underlying target ring is the
actual $A_{\mathrm{inf}}$; its identification $\mathrm{ainfId}$ in the
source merely records that ring, so the pure-tensor identity is
the one required there, including that identification.
\end{proof}
% END WPR22

% BEGIN WPR23
\paragraph{WPR23: all original Galois prism morphisms.}
\label{claim:WPR23}
For every $g\in G_K$ the original action
$\gamma_{\chi(g)}$ defines a prism automorphism
$gH(g):P_{\mathfrak S}\to P_{\mathfrak S}$, and
\[
 \mathrm{ioH}\circ gH(g)=\mathrm{ainfHom}(K,g)\circ\mathrm{ioH},
 \qquad \mathrm{ev}\circ gH(g)=\mathrm{ev}.
 \tag{WPR.28}
\]
These maps respect identity and multiplication in $G_K$.

\begin{proof}
The coefficient map $\sigma$ fixes $\mathbb Z_p$, so it commutes
with $\gamma_c$ for every $c\in\mathbb Z_p$. WPR3 gives
$\gamma_c\gamma_p=\gamma_p\gamma_c$. Thus
$\gamma_c\varphi_{\mathfrak S}=\varphi_{\mathfrak S}\gamma_c$;
WPR18 turns this into delta compatibility. For the unit
$c=\chi(g)$, WPR7 gives
\[
              \gamma_c Q_s=Q_s(1+U_s b_c).
\]
The second factor is a unit, because $U_s(0)=0$. The source
ideal is therefore preserved, in both directions by using
$c^{-1}$.

By WPR16 and the original theta equivariance,
$f\gamma_{\chi(g)}=\sigma_O(g)f$. The automorphism $g\in G_K$
fixes the actual embedded $\mathcal O_K$. The structure-class
criterion in WPR17 consequently proves base compatibility.
This establishes the claimed prism morphism and, using
$\gamma_{c^{-1}}$, its inverse. WPR16 proves the first equality
in (WPR.28) on underlying rings; prism morphisms with equal
underlying ring maps are equal. The second follows from
$\beta_c(0)=0$, so substitution preserves constant coefficients.
Finally WPR3 and the character law of WPR2 give identity and
composition exactly. No chosen lift of an element of $\Gamma$
enters this construction; it is defined for every original
$g\in G_K$.
\end{proof}
% END WPR23

% BEGIN WPR24
\paragraph{WPR24: the original scalar, parameter and distinguished polynomial maps.}
\label{claim:WPR24}
Put $\pi=T$, $Q=Q_s$, and let
$z:\mathbb Z_p[[T]]\to\mathfrak S$ be coefficient extension.
Then the original coefficient map $a$ satisfies
\[
 a\pi=\mathrm{let}\pi,\qquad aQ=R.Q,\qquad
 a\circ z=\mathrm{zlA},
 \tag{WPR.29}
\]
as literal equalities on the original rings. It is injective,
and it intertwines the prescribed Frobenius and Galois operators.

\begin{proof}
Injectivity and both intertwining identities are WPR16.
Each coefficient embedding sends the variable $T$ to $T$,
giving the first equality. Since every $e_i$ is
$\mathbb Q_p$-linear, its integral restriction commutes with
$\mathbb Z_p$ scalars. Coefficientwise this says precisely
$a\circ z=\mathrm{zlA}$; it is not a replacement scalar action.

For the polynomial equality, WPR6 identifies the supplied
$R.\varphi_0$ with substitution by $(1+T)^p-1$. Its defining
equation is $R.\varphi_0(T)=Tq_0$, so regularity of $T$ gives
\[
                     q_0=\sum_{j=0}^{p-1}(1+T)^j.
\]
The source's $R.Q$ is the $(s-1)$-fold full Frobenius iterate
of the constant tuple $q_0$. The tuple shift has no effect on
a constant tuple, and all the coefficient embeddings fix its
integer coefficients. WPR3 identifies its iterate with
\[
 \gamma_p^{s-1}(q_0)
   =\sum_{j=0}^{p-1}(1+T)^{j p^{s-1}}=Q_s.
\]
This proves $aQ_s=R.Q$ on every component. All cancellations
used regular $T$, not a division by $p$ or a change of layer.
\end{proof}
% END WPR24

\subsection{Both coefficient-coordinate assertions}

% BEGIN WPR25
\paragraph{WPR25: the full generic unit-separated coordinate theorem.}
\label{claim:WPR25}
Let $R_0$ be a nonzero commutative ring, $W$ a finite free
commutative $R_0$-algebra of rank $|J|$, $J$ a finite set,
and $O'$ a commutative $R_0$-algebra. Suppose maps
$\phi_k:W\to O'$, $k\in J$, are $R_0$-algebra maps and
for each $k\ne j$ there is $w\in W$ such that
$\phi_k(w)-\phi_j(w)$ is a unit. For every finite
$R_0$-basis $(o_i)_{i<d}$ of $O'$ and every $y\in(O')^J$
there is a unique $(c_i)_{i<d}\in W^d$ with
\[
                    y_k=\sum_{i<d}o_i\phi_k(c_i)
                    \quad(k\in J).                         \tag{WPR.30}
\]
This is the full original \texttt{existsUnique\_of\_separated}.

\begin{proof}
Consider the $O'$-algebra map
\[
 \Phi:O'\otimes_{R_0}W\longrightarrow(O')^J,
 \qquad \Phi(b\otimes w)_k=b\phi_k(w).
\]
Fix $k$. For every $j\ne k$, choose $w_{kj}$ with unit
$u_{kj}=\phi_k(w_{kj})-\phi_j(w_{kj})$. The element
\[
 t_k=\prod_{j\ne k}u_{kj}^{-1}
        \bigl(1\otimes w_{kj}-\phi_j(w_{kj})\otimes1\bigr)
\]
maps to the coordinate vector with value one at $k$ and zero
elsewhere: the $k$-coordinate of every factor is one, while
at coordinate $j$ its own factor is zero. Thus $\Phi$ is
surjective, since these vectors generate $(O')^J$ over $O'$.

Choose an $R_0$-basis of $W$ of size $|J|$. It makes source
and target of $\Phi$ free $O'$-modules of the same finite rank.
Choose preimages of the target basis to obtain a right-inverse
matrix $B$ for its matrix $M$. Then
$\det(M)\det(B)=1$, so the adjugate formula makes $M$ invertible.
This works over an arbitrary commutative $O'$; it does not
require a Noetherian hypothesis. Empty ranks satisfy the same
claim, or may be checked directly.

The given basis $(o_i)$ identifies $W^d$ with
$O'\otimes_{R_0}W$ by $(c_i)\mapsto\sum_i o_i\otimes c_i$.
Its composite with $\Phi$ is exactly the map in (WPR.30).
It is bijective, proving existence and uniqueness for every
supplied basis, with no preferred choice of that basis.
\end{proof}
% END WPR25

% BEGIN WPR26
\paragraph{WPR26: separators and the actual unramified coefficient coordinates.}
\label{claim:WPR26}
Let $K,LL$ be any original finite extensions and let
$e:\mathbb Z/n\mathbb Z\to\operatorname{Emb}_{\mathbb Q_p}(K_0,LL)$
be any injective enumeration, where $n=[K_0:\mathbb Q_p]$.
No $\mathrm{Good}$ or $\Gamma$ datum, or cyclic Frobenius
indexing, is assumed in this claim. These embeddings satisfy
the separation hypothesis of WPR25 on $W_0$. Consequently,
for every $\mathbb Z_p$-basis
$(o_i)_{i<d}$ of the original $O$ and every $y\in O^{\mathbb Z/n\mathbb Z}$,
there is a unique $c\in W_0^d$ with
\[
                    y_k=\sum_{i<d}o_i e_k(c_i).
 \tag{WPR.31}
\]
This proves the full original
\texttt{existsUnique\_unramifiedIntegers\_coords\_basis}.

\begin{proof}
First, in any local domain $D$, distinct $m$-th roots of unity
$a,b$ have unit difference whenever $m>0$ is a unit in $D$.
Indeed
\[
 (a-b)\sum_{j=0}^{m-1}a^j b^{m-1-j}=a^m-b^m=0.
\]
The nonzero first factor and the domain property make the sum
zero. If $a-b$ were a nonunit, their residues would agree.
Reducing the zero sum would give
$0=m\bar b^{m-1}$, a unit in the residue field, a contradiction.
This proves the source's generic unit-difference lemma.

Two distinct $\mathbb Q_p$-embeddings of $K_0$ must differ on
one of the prime-to-$p$ roots spanning $W_0$ in WPR1. Otherwise
they agree on $W_0$ and, after multiplying any field element
by a sufficiently large power of $p$, on all of $K_0$.
For such a root $\omega$ of order $m$ prime to $p$, both
$e_k(\omega)$ and $e_j(\omega)$ are integral roots of unity
in $O$, and are distinct. The number $m$ is a unit in
$\mathbb Z_p$ and hence in $O$. WIS1 gives that $O$ is a
local domain, so the preceding argument makes their difference
a unit. This is separation on the actual coefficient maps.

The supplied enumeration $e_k$ is injective. Its cardinality
is $n=[K_0:\mathbb Q_p]$, which equals the residue degree
and the size of the $\mathbb Z_p$-basis of $W_0$ in WPR1.
Apply WPR25 with $R_0=\mathbb Z_p$, $W=W_0$, $O'=O$.
All hypotheses use the original unramified field and the
supplied injective enumeration; no compatibility with Frobenius
was used. The result holds for every $d$ and every given
basis $(o_i)$, exactly as asserted. In particular it applies
to the embeddings of the original $\Gamma$ in the prism
construction.
\end{proof}
% END WPR26

% BEGIN WPR27
\paragraph{WPR27: the full coefficientwise formal-series lifting theorem.}
\label{claim:WPR27}
Let $A,B$ be commutative rings, $(o_i)_{i<d}$ a finite family
in $B$, and $e_k:A\to B$ ring maps indexed by an arbitrary
set $J$. Suppose every $y\in B^J$ has a unique $c\in A^d$
such that $y_k=\sum_i o_i e_k(c_i)$. Then every
$x\in B[[T]]^J$ has a unique $c\in A[[T]]^d$ such that
\[
                  x_k=\sum_{i<d}C(o_i)e_k(c_i).
 \tag{WPR.32}
\]
If $B$ is an $A$-algebra with finite basis $(b_i)$, this also
gives the full original \texttt{existsUnique\_\allowbreak of\_\allowbreak basis}:
\[
 x=\sum_{i<d}\operatorname{map}(A\to B)(c_i)C(b_i)
 \quad\text{with unique }c_i\in A[[T]].                      \tag{WPR.33}
\]

\begin{proof}
For each degree $r\geq0$ apply the hypothesis to the tuple
$([T^r]x_k)_{k\in J}$. Let its unique coordinate tuple be
$(c_{i,r})_{i<d}$ and define $c_i=\sum_{r\geq0}c_{i,r}T^r$.
The coefficient of degree $r$ on the right of (WPR.32) is
$\sum_i o_i e_k(c_{i,r})$, exactly $[T^r]x_k$. Thus the
series agree. Conversely an equality of series gives this
coefficient equation at every degree; the coefficient-level
uniqueness determines every $c_{i,r}$ and hence every $c_i$.
Only the finite sum over $i$ is involved; $J$ need not be
finite, and no topology on $A$ or $B$ has been introduced.

For (WPR.33), take $J$ to be a singleton, $e$ the algebra map,
and $o_i=b_i$. The coefficient-level hypothesis is exactly
the basis-coordinate theorem in $B$. Commutativity puts its
product in the displayed order. This proves the two original
generic formal-series statements at their full scopes.
\end{proof}
% END WPR27

% BEGIN WPR28
\paragraph{WPR28: both finite coordinate assertions on the actual rings.}
\label{claim:WPR28}
There are $d_0$ and $w_1,\ldots,w_{d_0}\in\mathfrak S$ such
that every $h\in\mathfrak S$ has a unique tuple
$c_i\in\mathbb Z_p[[T]]$ with
\[
                         h=\sum_i z(c_i)w_i.                \tag{WPR.34}
\]
In addition, for \emph{every} $d$, every $\mathbb Z_p$-basis
$(o_i)_{i<d}$ of $O$ and every original $x\in\mathrm{LetA}$,
there is a unique $c\in\mathfrak S^d$ such that
\[
                         x_k=\sum_{i<d}C(o_i)a(c_i)_k
                         \quad\text{for every }k.           \tag{WPR.35}
\]

\begin{proof}
For (WPR.34), use the $\mathbb Z_p$-basis
$\omega_1,\ldots,\omega_n$ of $W_0$ from WPR1 and set
$w_i=C(\omega_i)$. WPR27 gives existence and uniqueness
coefficient by coefficient. Its scalar map is precisely the
$z$ of WPR24, not an abstract finite-free identification.
Thus $d_0=n$ is one valid choice in the original existential
coordinate clause.

For (WPR.35), fix any of the basis data quantified in the
statement. WPR26 gives the coefficient-level unique
decomposition using these same $o_i$ and the original $e_k$.
Apply WPR27 with $A=W_0$ and $B=O$. By definition its series
maps $e_k$ are exactly the $k$-components of $a$. The result
is (WPR.35) for all of $\mathrm{LetA}$. No choice of the
$\mathbb Z_p$-basis of $O$ has been restricted, and no series
coefficient is inverted.
\end{proof}
% END WPR28

\subsection{Return to the two full original endpoints}

% BEGIN WPR29
\paragraph{WPR29: the full original Wach-prism existence theorem.}
\label{claim:WPR29}
For every original prime, $P,AI,K,s,LL$ with
$\mathrm{Good}(K,s,LL)$, every supplied $\Gamma,R,I$ and
$\mathrm{IoW}(I)$, the original
\texttt{exists\_\allowbreak wach\_\allowbreak prism\_\allowbreak over\_\allowbreak integers}
holds, with all sixteen output fields and no extra premise.
Specifically take $P_{\mathfrak S}$, $P_W$, $\mathrm{ev}$,
$\mathrm{ioH}$, $gH$, $a$, $z$, $\pi=T$ and $Q=Q_s$ above.
Their complete output is as follows:
\begin{enumerate}
\item $a$ is injective.
\item $I.\mathrm{io}(a h)=1\otimes\mathrm{ainfId}(\mathrm{ioH}(h))$
for every $h$.
\item $a(\varphi_{P_{\mathfrak S}}h)=\mathrm{let}\varphi_R(a h)$
for every $h$.
\item $a(gH(g)h)=\mathrm{let}\gamma_{R,\Gamma.\mathrm{res}(g)}(a h)$
for every $g,h$.
\item $\mathrm{ioH}\circ gH(g)=\mathrm{ainfHom}(K,g)\circ\mathrm{ioH}$
for every $g$.
\item $\mathrm{ev}\circ gH(g)=\mathrm{ev}$ for every $g$.
\item $a\pi=\mathrm{let}\pi$.
\item $\mathrm{ev}$ is surjective.
\item $\ker\mathrm{ev}=(\pi)$.
\item $P_{\mathfrak S}.I=(Q)$.
\item $aQ=R.Q$.
\item $a\circ z=\mathrm{zlA}$.
\item $P_{\mathfrak S}.A$ is a local ring.
\item $P_{\mathfrak S}.A$ is a domain.
\item The finite unique coordinates in (WPR.34) hold over this $z$.
\item For every original coefficient basis, the unique coordinates
in (WPR.35) hold over this $a$.
\end{enumerate}

\begin{proof}
WPR20 and WPR21 give the two bounded prisms over the same
$\mathcal O_K$, including their delta structures, invertible
ideals, mixed completeness, prism containment, bounded torsion
and structure maps. WPR21--WPR23 give all three kinds of
prism morphisms on these very rings. Thus the displayed
objects have the types required by the original endpoint.

Fields 1, 3 and 4 are WPR16. Field 2 is WPR15 and WPR22.
Fields 5 and 6 are WPR23; fields 7, 11 and 12 are WPR24.
Fields 8 and 9 are WPR21, fields 10, 13 and 14 are WPR20,
and fields 15 and 16 are WPR28. Each uses the same maps
$a,z,\mathrm{io}$ defined above. The original source's
$\mathrm{Good}$ data supplied precisely $s\geq1$ and
$K\subset K_s$ where those facts were needed in WPR17.
All primes, including $p=2$, and all permitted layer and
coefficient-field choices have been retained. This proves
both the source's \texttt{PrismS.exists\_prism} and its
exported statement in their original scope.
\end{proof}
% END WPR29

% BEGIN WPR30
\paragraph{WPR30: the full original conditional all-layer Wach existence.}
\label{claim:WPR30}
Retain the original
\texttt{Bhatt\allowbreak Scholze\allowbreak Crystalline\allowbreak Lattices}
input and the original $\mathrm{LocTied}(p,P,AI)$ premise.
Then the full original $\mathrm{WachEx}(AI)$ holds. Equivalently,
the theorem \texttt{WachLattice.main} is now reconstructed
at all of its original hypotheses. There is no remaining
period-ring supplier in this implication.

\begin{proof}
Fix an arbitrary original $\mathrm{Good}(K,s,LL)$ and arbitrary
permitted $\Gamma,R,I$, $\mathrm{IoW}(I)$ and supplied \'etale
module $\mathcal D$. We supply the three inputs isolated by
WLD26 on exactly these data.

First, WPR29 supplies the actual Wach prism, its maps and
both coordinate statements, including $a\circ z=\mathrm{zlA}$.
Second, WIS18 supplies injectivity and reflection of units
for the original $I.\mathrm{ioE}$. Third, IFD25 supplies the
full inhomogeneous fixed-Frobenius matrix theorem; taking its
right-hand side to be zero gives the homogeneous rectangular
descent required by WLD26. Its fixed subgroup is the original
$H=G_K\cap\ker\chi$. In particular this is the actual matrix
input, with every row and column dimension allowed; it is not
merely a scalar substitute.

For clarity, the resulting construction works on each original
crystalline lattice $T$. The retained crystalline-realization
input gives its realizing crystal on the actual Fontaine prism.
Evaluation on $P_{\mathfrak S}$ and $P_W$ supplies a finite
free lattice with Frobenius denominators controlled by $Q_s$.
The original comparison maps put its basis in the tensor period
module. Since $H$ has trivial image in $\Gamma$ (apply the
injective $\chi_\Gamma$ to the defining character identity),
this basis satisfies the fixed-matrix descent hypothesis. The
homogeneous IFD25 case places it in the original \'etale
coefficient module. WLD25 and WLD26 then prove every original
Wach-lattice field: the prescribed spanning equality, the two
Frobenius inclusions after powers of $Q_s$, the Galois
congruence modulo $T$, preservation by every allowed lattice
morphism, and the original trivial-object normalization. The
fullness clause of the retained crystalline-realization input
is the one used for the morphism and trivial-object returns;
the tied-period premise gives its original standard-period
identification.

The result of that construction is $\mathrm{WN}(\mathcal D,a\circ z)$.
WPR24 identifies its scalar map with the original $\mathrm{zlA}$,
so it is literally the required $\mathrm{NS}(\mathcal D)$.
Every layer and every supplied \'etale module were fixed only
after the universal inputs, and were arbitrary. This gives
the full $\mathrm{WachEx}(AI)$ and proves the stated original
implication.
\end{proof}
% END WPR30

\paragraph{Completed return and the next mathematical boundary.}
The prescribed prism and the full original conditional
\texttt{WachEx} theorem are now supplied. The latter retains
exactly its original tied-period and crystalline-realization
premises. The independent \texttt{FaithH} and \texttt{EtBasis}
producers, the Herr cohomology class maps, the quantitative
de Rham comparison and the common \texttt{ConclF} family still
lie ahead. Those are needed for the full \texttt{HW/Fon/Cmp}
data and the local row bounds. No whole
\texttt{GLOBAL\_\allowbreak DENOMINATOR} obligation is closed
by this chapter.

\paragraph{Source correspondence.}
All sources are at the accepted commit.
\par\noindent\texttt{7cbf7a0c32fca1badf05931879a37c3f8a2b03b4}.
Exact file, declaration, range and proof-block hashes accompany
each claim. The principal sources under
\texttt{Proofs/\allowbreak Inertia\allowbreak Eigenvalues} are
\texttt{Exists\allowbreak Wach\allowbreak Prism\allowbreak Over\allowbreak Integers},
\texttt{Exists\allowbreak Wach\allowbreak Inclusion\allowbreak Power\allowbreak Series\allowbreak Into\allowbreak One\allowbreak Tmul\allowbreak Ainf},
\texttt{Exists\allowbreak Unique\allowbreak Unramified\allowbreak Integers\allowbreak Coords\allowbreak Basis},
\texttt{Cyclotomic\allowbreak Integers\allowbreak Eq\allowbreak Witt\allowbreak Zeta\allowbreak Sub\allowbreak One},
\texttt{CDvd\allowbreak Of\allowbreak Map\allowbreak Eq\allowbreak Zero},
\texttt{Geom\allowbreak Sum\allowbreak One\allowbreak Add\allowbreak Prime},
and \texttt{Wach\allowbreak Lattices\allowbreak Exist\allowbreak All\allowbreak Layers\allowbreak Of\allowbreak Bhatt\allowbreak Scholze}.
The bounded-prism definitions and actual Fontaine object are
read on their original rings; the completed source theorem is
used as a contract for reconstruction, not as a substitute for
the human arguments above.
